% Worked Examples: Concept 3.7.4 - Integrator Windup
\documentclass[11pt,a4paper]{article}
\usepackage{worked-examples-template}

\title{\textbf{Worked Examples}\\[0.5em]
\large Concept 3.7.4: Integrator Windup\\[0.3em]
\normalsize \coursesubtitle{} -- \coursetitle}
\author{Assoc.\ Prof.\ William Robertson \and Dr Sean McGowan}
\date{\today}

\begin{document}

\maketitle
\tableofcontents
\newpage

\section{Concept 3.7.4: Anti-Windup}

\subsection{Example 1: Integrator Windup Problem and Solution}

\paragraph{Problem:} A PI controller is used to control a motor with position output. The controller is:
\begin{align}
u(t) = K_p e(t) + K_i \int_0^t e(\tau) d\tau
\end{align}

with $K_p = 5$ and $K_i = 2$. The control signal is limited to $|u| \leq 10$ volts.

(a) Explain what happens when a large step reference is applied that causes $u(t)$ to saturate.

(b) Design an anti-windup scheme using back-calculation with time constant $T_t = 1$ second.

\paragraph{Solution:}

\textbf{Part (a): Windup Problem}

When a large step reference is applied:

\begin{enumerate}
\item The error $e(t) = r(t) - y(t)$ is initially large
\item The controller output $u(t)$ quickly exceeds the limit and saturates at $u = 10$ V
\item While $u$ is saturated, the error remains large (output hasn't caught up to reference)
\item The integral term continues to accumulate: $\int e(\tau) d\tau$ keeps growing
\item This is called \textbf{integrator windup}
\item When the output finally approaches the reference, the integral has become very large
\item The controller output stays saturated even after the error becomes negative
\item This causes large overshoot and poor transient response
\end{enumerate}

\textbf{Part (b): Anti-Windup with Back-Calculation}

The anti-windup scheme adds a feedback path that prevents the integrator from winding up:

\begin{align}
\frac{dx}{dt} = K_i e - \frac{1}{T_t}(u_{sat} - u)
\end{align}

where:
\begin{itemize}
\item $x$ is the integral state
\item $u = K_p e + x$ is the controller output before saturation
\item $u_{sat} = \text{sat}(u)$ is the saturated control signal
\item $T_t = 1$ s is the tracking time constant
\end{itemize}

The saturated control is:
\begin{align}
u_{sat} = \begin{cases}
10 & \text{if } u > 10 \\
u & \text{if } -10 \leq u \leq 10 \\
-10 & \text{if } u < -10
\end{cases}
\end{align}

\textbf{How it works:}

\begin{itemize}
\item When $u$ is within limits: $u_{sat} = u$, so $\frac{dx}{dt} = K_i e$ (normal operation)
\item When $u$ saturates: The difference $(u_{sat} - u)$ is non-zero
\item This difference reduces the integration rate: $\frac{dx}{dt} = K_i e - \frac{1}{T_t}(u_{sat} - u)$
\item The integrator state is "tracked back" to a value consistent with the saturation limit
\item The time constant $T_t$ determines how quickly the tracking occurs
\item With $T_t = 1$ s, the integrator adjusts relatively quickly
\end{itemize}

\textbf{Implementation:}

The complete controller with anti-windup is:
\begin{align}
u &= K_p e + x \\
u_{sat} &= \text{sat}(u, -10, 10) \\
\frac{dx}{dt} &= K_i e - \frac{1}{1}(u_{sat} - u) = 2e - (u_{sat} - u)
\end{align}

\textbf{Benefits:}
\begin{itemize}
\item Prevents large overshoot caused by integrator windup
\item Maintains good tracking during saturation
\item Minimal impact on performance when not saturating
\item Simple to implement in both analog and digital systems
\end{itemize}
\end{document}
